\documentclass[11pt,a4paper]{article}
\usepackage[italian]{babel}
\usepackage[margin=2.5cm]{geometry}
\usepackage{parskip}
\usepackage{amsmath, amssymb, amsthm}
\usepackage{booktabs}
\usepackage{array}
\usepackage{xcolor}
\usepackage{needspace}
\usepackage{enumitem}
\usepackage{listings}
\usepackage{tikz}
\usepackage[most]{tcolorbox}
\usepackage[hidelinks]{hyperref}
\usetikzlibrary{decorations.pathreplacing, shapes.geometric, arrows.meta, positioning}

% Niente vedove/orfane: i paragrafi non lasciano righe isolate a fine/inizio pagina
\widowpenalty=10000
\clubpenalty=10000

% Elenchi più compatti
\setlist{itemsep=2pt, parsep=0pt, topsep=4pt, partopsep=0pt}

% --- Riquadri con bordo nero, senza numero (si spezzano tra due pagine) ---
% Uso: \begin{definizione} ... \end{definizione}
%      \begin{teorema}[Nome] ... \end{teorema}  (il nome è facoltativo)
\tcbset{nota/.style={enhanced jigsaw, breakable, colback=white,
  colframe=black, boxrule=0.5pt, sharp corners}}
\NewTColorBox{definizione}{}{nota,
  before upper={\textbf{Definizione.}\ }}
\NewTColorBox{teorema}{o}{nota,
  before upper={\textbf{Teorema\IfValueT{#1}{ (#1)}.}\ }}
\NewTColorBox{proposizione}{o}{nota,
  before upper={\textbf{Proposizione\IfValueT{#1}{ (#1)}.}\ }}
\NewTColorBox{proprieta}{o}{nota,
  before upper={\textbf{Proprietà\IfValueT{#1}{ (#1)}.}\ }}
\NewTColorBox{assioma}{o}{nota, boxrule=1pt,
  before upper={\textbf{Assioma\IfValueT{#1}{ (#1)}.}\ }}

% --- Ambienti semplici (senza riquadro) ---
% Il titolo sta in cima al blocco, anche se il contenuto inizia con un riquadro o
% con codice; e non resta mai da solo in fondo a una pagina.
% Uso: \begin{esempio}[Titolo facoltativo] ... \end{esempio}
\newcommand{\newrunin}[2]{%
  \NewDocumentEnvironment{#1}{o}{%
    \par\Needspace{4\baselineskip}\addvspace{\medskipamount}\noindent
    \textbf{#2}\IfValueT{##1}{ (##1)}\textbf{.}\ \ignorespaces}%
  {\par\addvspace{\medskipamount}}}
\newrunin{esempio}{Esempio}
\newrunin{esercizio}{Esercizio}
\NewTColorBox{osservazione}{}{nota,
  before upper={\textbf{Osservazione.}\ }}

% V e F nelle tavole di verità
\newcommand{\V}{\textsf{V}}
\newcommand{\F}{\textsf{F}}

% --- Codice C ---
% Uso: \begin{codice} ... \end{codice}      codice C numerato
%      \begin{comando} ... \end{comando}    comandi da terminale
%      \begin{risultato} ... \end{risultato} output di un programma
%      \lstinline|printf("ciao");|           codice dentro al testo
\lstdefinestyle{c}{
  language=C,
  basicstyle=\ttfamily\small,
  keywordstyle=\color{blue}\bfseries,
  commentstyle=\color{gray}\itshape,
  stringstyle=\color{red!70!black},
  numbers=left,
  numberstyle=\tiny\color{gray},
  numbersep=8pt,
  columns=flexible,
  keepspaces=true,
  breaklines=true,
  showstringspaces=false,
  tabsize=4,
  morekeywords={include, define},
  % lettere accentate nei commenti
  literate={à}{{\`a}}1 {è}{{\`e}}1 {é}{{\'e}}1 {ì}{{\`i}}1
           {ò}{{\`o}}1 {ù}{{\`u}}1 {È}{{\`E}}1 {À}{{\`A}}1
}
\lstset{style=c}
\newtcblisting{codice}{listing only, nota, left=6mm,
  listing options={style=c}}
\newtcblisting{comando}{listing only, nota,
  listing options={basicstyle=\ttfamily\small, numbers=none, language={},
    columns=flexible, keepspaces=true, breaklines=true}}
\newtcblisting{risultato}{listing only, enhanced jigsaw, breakable,
  colback=black!5, colframe=black, boxrule=0.5pt, sharp corners,
  listing options={basicstyle=\ttfamily\small, numbers=none, language={},
    columns=flexible, keepspaces=true, breaklines=true}}

% --- Diagrammi di flusso (simboli standard) ---
\tikzset{
  terminale/.style = {ellipse, draw, minimum width=2.5cm, minimum height=0.9cm},
  io/.style        = {trapezium, trapezium left angle=70, trapezium right angle=110,
                      draw, minimum width=2.5cm, minimum height=0.9cm},
  processo/.style  = {rectangle, draw, minimum width=2.5cm, minimum height=0.9cm},
  decisione/.style = {diamond, draw, aspect=2, minimum width=2.5cm},
  freccia/.style   = {thick, -{Stealth}}
}

\title{Il linguaggio C}
\author{Marco Cerbella}
\date{Lezione 3 -- 28 settembre 2026}

\begin{document}
\maketitle

\section{Linguaggi di programmazione}

\begin{definizione}
Un \emph{linguaggio di programmazione} è un linguaggio formale che specifica un
insieme di istruzioni per un elaboratore. Si usa per realizzare programmi che
implementano algoritmi. I linguaggi \emph{ad alto livello} sono vicini al
linguaggio umano e lontani dal linguaggio macchina.
\end{definizione}

I \emph{paradigmi di programmazione} classificano i linguaggi in base alle loro
caratteristiche. Il C è \emph{imperativo}: il programmatore indica alla
macchina \emph{come} modificare il proprio stato, ad esempio
\lstinline|a = a + 1;|.

\section{Dal linguaggio macchina al C}

\begin{itemize}
    \item Il \emph{linguaggio macchina} (sequenze di bit) è l'unico eseguito
          direttamente dal processore.
    \item L'\emph{assembly} usa abbreviazioni mnemoniche per le operazioni,
          ma è legato alla macchina.
    \item I linguaggi \emph{ad alto livello} sono indipendenti dalla macchina.
\end{itemize}

\begin{definizione}
Un \emph{compilatore} è un programma di sistema che traduce il programma
sorgente scritto in un linguaggio ad alto livello in un programma
\emph{target} equivalente (assembly/linguaggio macchina). Il programma target
viene eseguito in un secondo momento.
\end{definizione}

In questo corso il compilatore è \textbf{gcc}.

\begin{center}
\begin{tikzpicture}[node distance=1.6cm]
  \node (src) [processo] {sorgente C};
  \node (gcc) [processo, right=of src] {gcc};
  \node (exe) [processo, right=of gcc] {programma eseguibile};
  \draw [freccia] (src) -- (gcc);
  \draw [freccia] (gcc) -- (exe);
\end{tikzpicture}
\end{center}

\section{Caratteristiche del C}

Il C è un linguaggio ad alto livello, \emph{procedurale} e di uso generale,
progettato (Dennis Ritchie, anni '70) per scrivere sistemi operativi con la
massima indipendenza dall'hardware.

\textbf{Pregi:}
\begin{itemize}
    \item veloce (compilato, vicino alla macchina) e portabile;
    \item linguaggio piccolo e maturo, con molti strumenti;
    \item accesso diretto alla memoria e alle funzionalità di basso livello.
\end{itemize}

\textbf{Difficoltà:}
\begin{itemize}
    \item si può facilmente fare danni con accessi diretti alla memoria e
          puntatori;
    \item la memoria va gestita a mano;
    \item il codice è più verboso rispetto a linguaggi di scripting come Python.
\end{itemize}

\begin{osservazione}
Il linguaggio C in sé non ha istruzioni di input/output né di gestione della
memoria dinamica: queste funzioni sono fornite dalla \emph{libreria standard}
(ad esempio \lstinline|stdio.h|).
\end{osservazione}

\section{Standard}

Il C è stato standardizzato più volte: C89/C90 (ANSI/ISO), C99, C11, C18. Con
gcc si può scegliere lo standard:
\begin{comando}
gcc file.c -std=c11
\end{comando}

\section{Compilati e interpretati}

Un linguaggio \emph{compilato} (come C) viene analizzato a fondo dal traduttore
prima dell'esecuzione, e quasi tutti i controlli semantici sono fatti a tempo di
compilazione. In un linguaggio \emph{interpretato} (come Python) l'interprete
resta presente durante l'esecuzione e la maggior parte dei controlli è fatta a
tempo di esecuzione.
\end{document}
